\infographicchapter{MR Spectroscopy}{mrs.png}
\label{ch:infographic-mrs}

\section{Overview}

Magnetic resonance spectroscopy (MRS) is a specialized, non-invasive add-on to a standard MRI scan that checks the natural chemical makeup of living tissue. While a regular MRI uses signals from water to draw detailed pictures of body structures, MRS quiets the strong water signal so it can measure tiny natural chemicals (metabolites) inside a targeted spot, displaying them as a line graph of peaks---much like a non-invasive ``chemical check.'' It relies on the exact same physical principle as laboratory nuclear magnetic resonance (NMR) spectroscopy---doctors drop the word ``nuclear'' in the clinic simply to avoid confusion, since MRS uses harmless magnets and radio waves with zero radioactivity.

MRS is most often used for the brain. It measures key brain chemicals such as N-acetylaspartate (NAA, a marker of healthy nerve cells), choline (Cho, which rises when cells are growing or turning over rapidly), creatine (Cr, a steady energy marker used for comparison), and lactate (Lac, a sign of low oxygen or energy stress). Doctors use these chemical patterns to help evaluate brain tumors, tell the difference between a returning tumor and harmless scar tissue from past radiation, distinguish a pus-filled brain infection (abscess) from a tumor, check low-oxygen brain injury, and evaluate inherited metabolic conditions.

Because MRS runs on the exact same MRI scanner, it adds valuable biochemical clues without X-rays, needles, or extra dye. Its main limits are that it measures a fairly large chunk of tissue at once (so very tiny spots are hard to isolate), it can be disrupted near skull bone, air sinuses, bleeding, or metal plates, and doctors usually interpret chemical ratios alongside standard MRI rather than on their own. Standard MRI safety rules for metal implants and claustrophobia apply.

\section{Underlying Physical Principles}
\begin{itemize}
\item \textbf{Chemical Signatures in a Magnetic Field:} Inside a strong, ultra-even MRI magnet (1.5 to 3.0~Tesla), hydrogen atoms inside different natural body chemicals experience slightly different local magnetic environments. This tiny frequency difference---called chemical shift and measured in parts per million (ppm)---lets the scanner tell one chemical from another.
\item \textbf{Quieting Water and Fat:} Because water is thousands of times more abundant in tissue than natural metabolites, the scanner first sends targeted radio pulses to quiet the strong water signal and block nearby scalp or skull fat so the tiny chemical peaks can be heard clearly.
\item \textbf{Zooming Into a Targeted Box (Voxel):} Intersecting radio pulses---typically using specialized localization methods called PRESS (Point Resolved Spectroscopy) or STEAM (Stimulated Echo Acquisition Mode)---isolate signals strictly from a carefully chosen small box (single-voxel MRS) or a checkerboard grid of boxes (multi-voxel MRS, also called chemical shift imaging or CSI) placed right over the area of interest. While routine hospital scans listen to hydrogen protons ($^1$H), specialized research scans can also tune into phosphorus ($^{31}$P) to study cellular energy.
\item \textbf{Timing Echoes to Separate Overlapping Peaks:} Adjusting how long the scanner waits before listening to the echo (short versus medium echo time) helps separate chemicals whose signals sit close together---for instance, flipping the lactate peak upside down at a medium echo time so it stands apart from fat.
\item \textbf{Building the Chemical Graph:} Computer software turns the captured signals into a smooth graph of chemical peaks and calculates helpful ratios---such as choline compared to NAA---for the specialist to interpret.
\end{itemize}

\section{Imaging Mechanism}
\begin{enumerate}
\item \textbf{Head Coil Setup and Standard MRI First:} You lie comfortably inside a padded head coil (a 3.0~Tesla scanner is often preferred for sharper chemical separation). The scanner first takes standard structural MRI pictures so the specialist can place the small measurement box accurately on the target spot and on healthy comparison tissue.
\item \textbf{Fine-Tuning the Magnet and Quieting Water:} The scanner automatically fine-tunes (``shims'') the magnetic field inside the chosen box to make it as smooth as possible, quiets the water and surrounding fat signals, and runs the spectroscopy pulses.
\item \textbf{Collecting Faint Chemical Signals:} Because natural chemical signals are very faint, the receiver coil listens and averages many quick repetitions together to build a clean, noise-free reading.
\item \textbf{Processing the Chemical Peaks:} Dedicated spectroscopy software smooths out background ripples, removes any leftover water signal, and maps each distinct chemical peak along the ppm scale.
\item \textbf{Side-by-Side Chemical and Image Review:} The final chemical line graphs and color-coded ratio maps are laid right next to your standard brain MRI slices so the neuroradiologist can evaluate both appearance and chemistry together.
\end{enumerate}

\section{Key Technical Parameters}
\noindent
\begin{tabularx}{\textwidth}{@{}l X@{}}
\toprule
\textbf{Parameter} & \textbf{Typical Value / Description} \\
\midrule
Physical Principle & MRI tuned to suppress water and measure tiny natural chemical differences (chemical shift in ppm) inside tissue \\
Energy / Signal Source & High-field static magnet (1.5--3.0~T; 3.0~T preferred for clearer peak separation), magnetic gradients, and radio pulses \\
Typical Frequency / Energy & Non-ionizing radio waves (64~MHz at 1.5~T; 128~MHz at 3.0~T) \\
Spatial Resolution & Box size of 1--8~mL (about 1--2~cm wide) for single-voxel scans, or 0.5--1.5~mL per box in multi-voxel grids \\
Typical Scan Duration & 10--20~minutes for the MRS portion; 40--60~minutes total when combined with a full brain MRI \\
Contrast Agent & None required for MRS itself (run before or independently of routine MRI contrast dye) \\
Ionizing Radiation & None (0~mSv; zero X-ray exposure) \\
\bottomrule
\end{tabularx}

\section{Common Symptoms That Prompt This Scan}
\begin{itemize}
\item \textbf{Unclear Brain Mass Found on Standard MRI:} Headaches, seizures, or neurological changes where standard MRI shows a spot and doctors need chemical clues to tell tumor, infection, and inflammation apart.
\item \textbf{New Symptoms After Brain Tumor Radiation:} Returning headaches or weakness after treatment, when MRS helps distinguish active tumor regrowth from radiation-related tissue injury.
\item \textbf{Developmental Delay or Unexplained Seizures in Children:} Delayed milestones, muscle weakness, or metabolic symptoms pointing to rare inherited brain chemistry disorders.
\item \textbf{Prostate or Specialized Organ Workups:} Elevated PSA or unclear standard MRI findings in select specialized centers evaluating tissue chemistry.
\end{itemize}

\section{Clinical Usages}

\subsection{Primary Usages}
\begin{itemize}
\item \textbf{Evaluating and Grading Brain Tumors:} Checks a brain mass for elevated choline (rapid cell growth) and dropped NAA (loss of normal nerve cells), helping tell tumors apart from non-tumor conditions and estimate how aggressive a glioma may be.
\item \textbf{Returning Tumor vs. Radiation Scar Tissue:} After brain radiation therapy, helps distinguish a truly returning tumor (high choline) from harmless radiation tissue injury (low choline, low creatine, and low NAA with dead-cell lipid/lactate peaks).
\item \textbf{Telling Brain Abscesses Apart from Tumors:} Ring-shaped spots on MRI can be either an infection (pus-filled abscess) or a tumor. Finding specific amino acid peaks from breaking-down white blood cells strongly points to a treatable abscess rather than cancer.
\item \textbf{Metabolic and Low-Oxygen Brain Conditions:} Identifies distinct chemical clues in inherited metabolic disorders (such as Canavan disease, creatine deficiency, or mitochondrial disorders) and low-oxygen brain injury in newborns.
\end{itemize}

\subsection{Secondary Usages}
\begin{itemize}
\item \textbf{Guiding Brain Biopsies:} Uses a multi-voxel grid to pinpoint the single most active region (highest choline-to-NAA ratio) inside an uneven tumor so surgeons sample the most informative spot (achieving a 100\% diagnostic tissue yield across 42~patients in one clinical series).
\item \textbf{Fine-Tuning Radiation Therapy Targets:} Uses 3D multi-voxel chemical grids to detect microscopic tumor activity extending beyond the edge visible on standard MRI (seen in 45\% of grade~II glioma patients in one study), helping radiation teams shape treatment fields more accurately.
\item \textbf{Locating Seizure Focuses:} Helps identify which side of the brain is driving temporal lobe epilepsy by spotting lowered NAA even when regular MRI pictures look nearly normal.
\item \textbf{Measuring Liver Fat:} Accurately measures liver fat levels in fatty liver disease, especially in specialized liver evaluations and clinical research.
\item \textbf{Selected Prostate or Body Checks:} Occasionally used as an extra tool in prostate or body evaluations, though modern diffusion and contrast MRI have largely taken over those roles.
\end{itemize}

\section{Advantages}
\begin{itemize}
\item \textbf{Non-Invasive Chemical Check:} Reveals natural chemical and metabolic changes inside living tissue without surgery or needles, giving information that standard pictures alone cannot show.
\item \textbf{Sharpens Uncertain Diagnoses:} Helps clarify tricky MRI findings---such as returning tumor versus radiation scarring, or infection versus tumor---sometimes sparing patients from unnecessary brain biopsies; in a large clinical study of unclear brain lesions, adding MRS raised diagnostic accuracy from 55\% with standard MRI alone to 71\%.
\item \textbf{Zero Radiation and No Extra Dye:} Uses no X-rays and requires no special contrast injection for the spectroscopy measurement itself.
\item \textbf{Fits Right Into a Standard MRI Visit:} Runs on the same scanner right alongside a routine brain MRI and can help target biopsies accurately if tissue sampling is needed.
\end{itemize}

\section{Disadvantages}
\begin{itemize}
\item \textbf{Large Measurement Boxes:} Because chemical signals are faint, the scanner needs a relatively large box (1 to 8~mL). For very small spots, the box may accidentally include normal tissue, fluid, or dead tissue around the edge, blurring the result.
\item \textbf{Disrupted Near Bone, Air, Blood, or Metal:} Areas close to the skull base, air-filled sinuses, scalp fat, bleeding, calcium, or metal surgical plates disturb the magnetic field and can make graphs unreadable.
\item \textbf{Overlapping Patterns:} Some non-cancerous inflammatory spots (like sudden demyelination) can temporarily mimic a tumor's chemical graph, and overlapping fat and lactate peaks require careful echo timing.
\item \textbf{Uses Ratios Rather Than Exact Amounts:} Doctors usually interpret relative peak ratios (like choline compared to NAA) rather than exact concentrations, and results depend heavily on clean setup and magnet fine-tuning.
\item \textbf{Adds Extra Time and Shares MRI Limits:} Extends the MRI visit by 10 to 20~minutes (making stillness even more important) and carries all standard MRI rules regarding metal implants and claustrophobia.
\end{itemize}

\section{Contraindications}

\subsection{Absolute}
\begin{itemize}
\item \textbf{Unsafe Metal Implants:} Non-MR-conditional pacemakers or defibrillators, magnetic brain aneurysm clips, cochlear implants, metal fragments in the eyes, and non-MR-conditional nerve stimulators.
\end{itemize}

\subsection{Relative}
\begin{itemize}
\item \textbf{Local Field Disruptors:} Heavy bleeding inside the target spot, dense calcium, metal skull plates, or spots tucked right next to the skull base or air sinuses can prevent clean fine-tuning and make graphs unreadable.
\item \textbf{Trouble Lying Perfectly Still:} Agitation, tremor, or severe claustrophobia that makes lying very still difficult (unless safe MRI-monitored sedation is arranged).
\end{itemize}

\section{Patient Safety Guidelines}
\begin{itemize}
\item \textbf{Thorough Metal Screening:} Everyone completes a full MRI safety questionnaire beforehand to verify safety for brain aneurysm clips, nerve stimulators, shunts, pacemakers, or metal fragments.
\item \textbf{Preventing Skin Warmth:} Especially on stronger 3.0~Tesla scanners, water-quieting pulses deposit a bit more radiofrequency energy. Staff keep scanner settings strictly within safety limits and use soft foam pads to keep skin away from tunnel walls and cables.
\item \textbf{Hearing Protection:} Fitted earplugs or MRI-safe headphones protect your ears from the loud tapping and humming sounds during both standard MRI and spectroscopy.
\item \textbf{Extra-Comfortable Head Support:} Because even tiny head movements spoil chemical box alignment, staff cushion your head snugly and comfortably, keep an intercom and squeeze-ball call button ready, and use MRI-safe monitoring if sedation is needed.
\item \textbf{Standard Dye Rules If Paired with Contrast MRI:} MRS itself needs no dye, but if your brain MRI includes IV gadolinium contrast, standard kidney and pregnancy safety checks apply.
\end{itemize}

\section{Prerequisite Tests}
\begin{itemize}
\item \textbf{No Routine Lab Work for MRS Alone:} Because MRS measures natural body chemicals without contrast dye, no blood tests are needed for spectroscopy itself (kidney checks apply only if paired with contrast MRI).
\item \textbf{Standard Structural MRI First:} Regular multi-angle brain MRI pictures are always acquired first so the team can see the exact spot, check for bleeding or calcium, and position the chemical measurement box accurately.
\item \textbf{MRI Safety Check and Implant Cards:} Complete the standard MRI safety form, have an eye X-ray if you have a history of metal eye injury, and bring device cards for any skull plates, shunts, clips, or implants.
\item \textbf{Prior Surgery, Radiation, or Metabolic Records:} Knowing dates of past brain surgery or radiation treatments---or reviewing blood and urine metabolic tests in inherited conditions---helps the neuroradiologist choose the best settings and interpret the graphs accurately.
\end{itemize}

\section{Preparing for the Scan}
\begin{itemize}
\item \textbf{Metal-Free Gown and Skipping Metallic Cosmetics:} Change into a hospital gown and remove jewelry, hairpins, glasses, hearing aids, removable dentures, and foil-backed skin patches. For brain MRS, skip metallic hair products or heavy eye makeup, which can disturb magnet fine-tuning.
\item \textbf{Bring Implant Cards and Past Records:} Bring device cards for any skull plates, shunts, clips, or pacemakers, plus any outside brain MRI or radiation records.
\item \textbf{Eating, Medicines, and IV Setup:} No fasting is needed for routine brain MRS unless sedation or a liver spectroscopy scan (4- to 6-hour fast) is planned. Take your usual medicines and expect a quick arm IV if paired with contrast brain MRI.
\item \textbf{Restroom Visit and Head Stillness:} Use the restroom beforehand since adding MRS makes the MRI 10 to 20~minutes longer. Wear the supplied ear protection and let staff place soft foam pads around your head to help you stay effortlessly still.
\item \textbf{Claustrophobia or Neck Comfort:} Let the clinic know ahead of time if claustrophobia or neck pain makes lying flat difficult so extra cushions or mild relaxing medicine (with a driver home) can be arranged.
\end{itemize}

\section{During Scan}
\begin{itemize}
\item \textbf{Positioning and Staying Extra Still:} You lie on your back with your head resting inside a cushioned head coil. Because MRS listens to tiny chemical signals from one exact spot, staying completely motionless---avoiding nodding, talking, or coughing---makes a big difference.
\item \textbf{Breathing Instructions:} Breathe calmly and naturally during brain MRS. If liver or belly spectroscopy is performed, the technologist will guide you through gentle breath-holds or steady quiet breathing.
\item \textbf{Sounds and Sensations:} MRS feels just like a regular MRI---completely painless and needle-free, with loud rhythmic clicking, humming, and buzzing softened by ear protection. You may notice gentle warmth during higher-field 3.0~Tesla water-quieting pulses.
\item \textbf{Staying in Touch:} You hold a squeeze-ball call button the entire time and can talk with the technologist between scan blocks. Speak up right away if you feel localized warmth, headache, discomfort, or anxiety.
\item \textbf{How Long It Takes:} The spectroscopy part adds about 10 to 20~minutes to a standard brain MRI, bringing the full visit to around 40 to 60~minutes.
\end{itemize}

\section{Normal Appearance}
\begin{itemize}
\item \textbf{Healthy Nerve Peak on Top:} In healthy adult brain tissue, the NAA peak (healthy nerve cells) is the tallest peak on the graph, followed step by step by creatine and then choline, forming a smooth upward slope (called Hunter's angle).
\item \textbf{Balanced Ratios:} Healthy brain choline-to-creatine and choline-to-NAA ratios stay below 1.0, showing steady cell stability without abnormal growth.
\item \textbf{Clean Flat Baseline:} Normal graphs show a smooth, flat bottom line with water nicely quieted and no big spikes of lactate (low oxygen), lipids (dead tissue), or abscess amino acids.
\end{itemize}

\section{Observing for Abnormal Patterns}
\begin{itemize}
\item \textbf{High Choline and Low NAA (Tumor Pattern):} A tall choline spike (rapid cell membrane growth) paired with a dropped NAA peak (loss of normal nerve cells) reverses the normal upward slope and strongly suggests an active or returning brain tumor.
\item \textbf{Lactate and Lipid Spikes (Low Oxygen or Dead Tissue):} A double-peaked lactate signal (which flips upside down at medium echo time) points to low oxygen in stroke, mitochondrial disease, or aggressive tumors, while broad lipid peaks reflect dead cells in high-grade tumors or radiation scar tissue.
\item \textbf{Amino Acid Peaks in Brain Abscesses:} Distinct spikes of amino acids (along with acetate and succinate) while normal brain peaks disappear help distinguish a pus-filled brain abscess from a dead-center tumor.
\item \textbf{Condition-Specific Clues:} Elevated myo-inositol can appear in gliosis, low-grade tumors, or Alzheimer disease; an unusually tall NAA peak is the hallmark of Canavan disease; and an alanine peak supports a benign meningioma.
\item \textbf{Harmless Technical Look-Alikes:} Placing the measurement box too close to skull bone, scalp fat, air sinuses, or old bleeding can create fake fat spikes or wide, messy peaks that mimic dead tumors. Sudden inflammatory spots (tumefactive demyelination) can also briefly imitate a tumor's graph.
\end{itemize}

\section{Dealing with Uncertainties}
\begin{itemize}
\item \textbf{Magnet Fine-Tuning Challenges:} Spots right next to the skull base, air sinuses, bleeding, calcium, or metal plates can distort the magnetic field and blur peaks. Technologists adjust fine-tuning manually, shift the box slightly away from bone or air, or use smaller multi-voxel grids.
\item \textbf{Box-Edge Mixing and Fat Spillover:} Because measurement boxes are fairly large, small spots can mix with surrounding healthy brain or fluid, and nearby scalp fat can spill in. Outer-edge blocking bands and careful box checks on standard MRI slices help keep readings pure.
\item \textbf{Overlapping Peaks and Mimics:} Fat and lactate sit close together on the graph, so running both short and medium echo times separates them cleanly. Because inflammation can sometimes mimic a tumor, doctors never judge a graph in isolation.
\item \textbf{Combining All Clues:} Neuroradiologists always read MRS alongside contrast MRI, diffusion scans, blood-flow (perfusion) MRI, and your medical history. If a spot stays uncertain, a repeat scan in 2 to 3~months or a targeted biopsy provides a definite answer.
\end{itemize}

\section{Augmenting Methods}
\begin{itemize}
\item \textbf{Blood-Flow (Perfusion) MRI:} Measuring blood volume inside a brain spot pairs closely with MRS---returning tumors usually have rich blood flow, while radiation scar tissue has low blood flow.
\item \textbf{Diffusion MRI and Fiber-Tract Mapping (DTI):} Diffusion scans confirm thick pus inside an abscess alongside MRS amino acid peaks, while fiber-tract mapping shows surgeons where important nerve pathways run around a tumor.
\item \textbf{Amino Acid PET Scans:} Specialized PET scans (such as FET or FDOPA PET) offer sharper detail near surgical beds or bleeding when MRS graphs are disturbed.
\item \textbf{Molecular Biopsy Testing:} If tissue is sampled, lab and genetic testing (such as IDH mutation status) gives the definitive tumor classification.
\item \textbf{Metabolic Blood/Urine Tests and Fitting Software:} Blood, urine, and spinal fluid panels confirm inherited metabolic disorders, while automated curve-fitting software makes peak ratios more consistent.
\end{itemize}

\section{Comparison with Alternative Modalities}
\begin{itemize}
\item \textbf{Compared with Standard Brain MRI:} Standard MRI shows exact size, shape, swelling, and dye uptake with sharp millimeter detail, while MRS adds chemical clues to tell tumors, radiation scarring, infections, and metabolic conditions apart.
\item \textbf{Compared with Blood-Flow (Perfusion) MRI:} Perfusion MRI is faster, sharper, and less disrupted by bone or bleeding, while MRS directly checks nerve health, cell turnover, and chemical markers like lactate or abscess amino acids.
\item \textbf{Compared with Amino Acid PET:} Amino acid PET gives sharper views near surgical scars or bleeding, whereas MRS uses zero radiation, needs no radioactive tracer, and fits right into a routine MRI visit.
\end{itemize}

\section{Follow-up Steps}
\begin{itemize}
\item \textbf{Immediate Return to Your Day:} Because MRS itself uses no dye, you can resume normal activities right after the scan (following routine IV checks or sedation recovery if those were part of your full MRI visit).
\item \textbf{Checking Graph Quality and Full MRI Together:} Your neuroradiologist checks that the peaks are crisp, water is cleanly quieted, and the measurement box sat squarely inside the target spot before combining results with your full MRI.
\item \textbf{Guiding Biopsies or Targeted Treatment:} If a biopsy is needed, multi-voxel chemical maps help surgeons aim for the most active part of a lesion. Suspected abscesses are paired with diffusion MRI for prompt treatment, and suspected metabolic disorders are confirmed with blood, urine, or genetic tests.
\item \textbf{Follow-Up Tracking After Radiation:} When watching an uncertain spot after radiation therapy, repeating MRS alongside standard and blood-flow MRI every 2 to 3~months shows whether peaks are fading (healing scar tissue) or choline is climbing (returning tumor).
\end{itemize}

\section{Key References}
\begin{itemize}
\item Oz G, Alger JR, Barker PB, et al. Clinical proton MR spectroscopy in central nervous system disorders. \textit{Radiology}. 2014;270(3):658--679.
\item American College of Radiology (ACR). \textit{ACR--ASNR--SPR Practice Parameter for the Performance and Interpretation of Magnetic Resonance Spectroscopy of the Central Nervous System}. Reston, VA: ACR.
\item American College of Radiology (ACR). \textit{ACR Manual on MR Safety}. Reston, VA: ACR; current edition.
\item Weinberg BD, Kuruva M, Shim H, Mullins ME. Clinical applications of magnetic resonance spectroscopy in brain tumors: from diagnosis to treatment. \textit{Radiologic Clinics of North America}. 2021;59(3):349--362.
\item Hollingworth W, Medina LS, Lenkinski RE, et al. \textit{A Systematic Literature Review of Magnetic Resonance Spectroscopy for the Characterization of Brain Tumors}. National Library of Medicine (NIH PubMed Central; AJNR Am J Neuroradiol, 2006). Available at: \url{https://pmc.ncbi.nlm.nih.gov/articles/PMC7977543/}.
\item Radiopaedia.org. \textit{MR spectroscopy}. Radiopaedia. Available at: \url{https://radiopaedia.org/articles/mr-spectroscopy-1}.
\end{itemize}

\section{Further Reading}
\begin{itemize}
  \item Wikipedia contributors. \textit{In vivo magnetic resonance spectroscopy}. Wikipedia, The Free Encyclopedia. Available at: \url{https://en.wikipedia.org/wiki/In_vivo_magnetic_resonance_spectroscopy}.
\end{itemize}
